\section{Evaluation design}
\label{sec:evaluation}

We fixed hypotheses, comparators, measures and statistics in an analysis plan committed before the first held-out model call \cite{casepath2026repo}. Appendix~\ref{app:deviations} lists every deviation from it.

\subsection{Hypotheses}
\label{sec:hypotheses}

The plan states five hypotheses, given here as registered.
\begin{description}[leftmargin=1.2em,labelindent=0pt]
\item[H1 Safety.] On the tenancy claims, \casepath with exact deferral makes no more unsafe errors than \design{Rules in context}.
\item[H2 Deferral.] On the tenancy claims, \casepath defers more required documents than \design{Rules in context}; exact deferral defers fewer than the production engine without more unsafe errors.
\item[H3 Changed facts.] When one fact of a theft case changes, \casepath makes fewer wrong changes than \design{Direct} and \design{Rules in context}.
\item[H4 Unresolved facts.] When one fact of a theft case is left open, \casepath makes fewer premature requests and other changes than \design{Direct} and \design{Rules in context}.
\item[H5 Stability.] \casepath's run-to-run agreement exceeds that of the model-written plans.
\end{description}
Each primary test contrasts \casepath with \design{Rules in context}, fixed in advance as the strongest comparator; contrasts with \design{Direct} and with majority voting are secondary. The \emph{primary analysis} uses \texttt{gpt-6.1-terra}, a frontier model. The \emph{secondary analysis}, registered as a model-sensitivity check, repeats \casepath and \design{Rules in context} with \texttt{gpt-6-luna}, the cost-efficient model that the deployed agent uses.

\subsection{Data: two regimes}
\label{sec:data}

The two datasets come from the original \casepath study \cite{casepath2026repo}, with one new set of variants (Table~\ref{tab:data}).

\paragraph{Clean rules, complete references (tenancy)}
150 synthetic legal-expenses claims about tenancy defects, lease terminations and rent increases arrive as an inbox receives them: a message in English or Swiss Standard German, by e-mail, portal, chat or phone note, with attachments in 45 claims. The rules are the three admitted process templates. Each claim's reference lists the documents it requires at intake, generated from the claim's scenario and the same templates the designs receive: with the scenario's true conditions, the \casepath engine reproduces the reference for 147 of 150 claims, and the other three differ only by a photograph that no text-only design can read, which we do not score. This circularity is deliberate: the regime measures how a design applies correct rules to incomplete claims, not whether the rules are legally right. The generator gives every fact a value. Facts a message never mentions are false in the scenario, and the two routing facts are always true, although some messages never state them. We evaluate on the 90 claims of 17 held-out scenarios; six development claims served for piloting.

\paragraph{Extracted rules, changed facts (theft)}
36 pairs of household-theft cases differ in one sentence that switches one of nine conditions from false to true, for instance from ``No stolen item is separately scheduled in the policy with an insured value of CHF 1,000 or more'' to ``A stolen watch is separately scheduled in the policy with an insured value of CHF 4,000''. Each pair's reference lists the checklist changes that the sources justify, with alternatives where either document suffices. The rules were extracted by a language model from 197 source passages and are imperfect, which is what this regime is for. One of four stories is development data; the other three give 27 held-out pairs.

\paragraph{Unresolved-fact variants (new)}
For six conditions whose documents are downstream requirements, we replaced the condition's sentence with one that leaves it open, for example ``A stolen watch may be separately scheduled in the policy, but the policyholder has not yet found the policy schedule to confirm this'' (Appendix~\ref{app:unresolved}). The correct response is to change no request and to hold the condition's documents until the fact is clarified. This yields 18 held-out pairs. We excluded the other three conditions because their document is itself the evidence that settles the fact.

\paragraph{Reference check}
To test whether the theft references depend on how they were built, a separate model run that saw only the source passages and the catalogue re-derived, three times per condition, the documents each condition change requires. For seven of nine conditions it reproduced the reference exactly in all three runs. For the formal expert procedure it also required the experts' findings; for the insurer's declination it required no document, where the reference requires the declination letter. We report results on all conditions and, as a sensitivity analysis, on the seven agreed ones.

\begin{table}[t]
\caption{Evaluation data. Counts refer to the held-out sets.}
\label{tab:data}
\small
\begin{tabularx}{\linewidth}{@{}lXX@{}}
\toprule
& Tenancy claims & Theft cases \\
\midrule
Regime & clean rules, complete references & extracted rules, changed or unresolved facts \\
Units & 90 claims from 17 scenarios & 27 changed-fact pairs, 18 unresolved-fact pairs (3 stories) \\
Input & message and attachment text; 2 languages, 4 channels & 141--155-word English case with nine condition statements \\
Rules & 3 admitted process templates, 9--11 document types each & 98 extracted rules over 87 conditions; 33 document types \\
Reference & required documents at intake (generator state) & required change per condition (source contract) \\
Hypotheses & H1, H2, H5 & H3, H4 \\
\bottomrule
\end{tabularx}
\end{table}

\subsection{Compared designs}
\label{sec:designs}

All designs receive the same case text, rules, document catalogue and case-reading instructions, use low reasoning effort and schema-constrained output, and differ only in who applies the rules and whether judgments are verified (Table~\ref{tab:designs}). \design{Direct} sees the theft source passages and the catalogue and writes the checklist. \design{Rules in context} additionally sees the extracted rules (theft) or the admitted templates (tenancy) and writes the evidence plan, with the option to hold documents. \design{Self-consistency} takes, for every document, the majority decision of the three runs of a model-written design. \design{Model executes rules} judges every condition and then applies the stated three-valued process semantics itself, in the same output, with exactly the information \casepath's code uses. \design{Two-valued execution} re-executes \casepath's recorded judgments with unresolved facts read as false or as true, which isolates the semantics; it was added after the runs. For \casepath we report the production engine (strong Kleene evaluation) and exact deferral (supervaluation), the latter also without the verification call; these variants reuse the same model outputs.

The grid is uneven for three reasons. \design{Direct} gives the model the legal sources instead of rules, and only the theft data include such sources. \design{Model executes rules} tests whether a model can carry out the process semantics, which matter only where requirements depend on process steps, as in tenancy. The secondary model ran, as planned, only \casepath and the strongest comparator.

\begin{table}[t]
\caption{Compared designs. ``Applies rules'' names who turns case facts into requests.}
\label{tab:designs}
\small
\begin{tabularx}{\linewidth}{@{}lXlll@{}}
\toprule
Design & Representative of & Applies rules & Verified & Data, models \\
\midrule
\design{Direct} & LLM with policy text in the prompt \cite{yao2025taubench} & model & no & theft; primary \\
\design{Rules in context} & process-aware prompting \cite{xiao2024flowbench} & model & no & both; both \\
\design{Self-consistency} & majority of three runs \cite{wang2023selfconsistency} & model & no & both; both \\
\design{Model executes rules} & stepwise rule application \cite{servantez2024chain} & model & no & tenancy; primary \\
\design{Two-valued execution} & extraction plus rule engine \cite{pan2023logic,omg2024dmn} & code & no & both; both (post hoc) \\
\casepath & this article & code & yes & both; both \\
\bottomrule
\end{tabularx}
\end{table}

\subsection{Measures}
\label{sec:measures}

On the tenancy claims, every required document has one of four outcomes: \emph{requested}, \emph{held} behind a question, \emph{handed over} with the whole claim to a person, or \emph{missed}. Every requested document that is not required is an \emph{unnecessary request}. Unsafe errors are unnecessary requests plus missed documents, reported per claim; deferrals are held plus handed-over documents. The plan did not say how to score handed-over documents. We decided after the runs not to count them as errors, because a person then decides, and report a \emph{strict} variant that counts them as missed wherever H1 is discussed. No planned measure penalises holding a document that is not required; we report such holds descriptively.

On a theft pair, the signed difference between the two checklists is scored against the reference change with the frozen scorer of the original study: a change is correct, wrong (an addition or withdrawal the sources do not justify) or missed, pooled into precision, recall and $F_1$. On an unresolved-fact pair, a response is \emph{premature} if it requests a document of the open condition, and \emph{stable} if it changes no request; H4 counts every change, premature or not. Run-to-run agreement is the mean pairwise Jaccard similarity of the request lists of the three runs of the same input. Cost is the provider charge.

\subsection{Protocol and statistics}
\label{sec:protocol}

Every design ran three times on every held-out case with each of its models, on 9 November 2026 through OpenRouter; responses, token counts and charges are logged. A failed output is scored as an empty plan: one \design{Direct} output failed, and the original scorer had dropped it, which we corrected. Units are claims and pairs, averaged over runs, and uncertainty is dominated by clusters: the 17 tenancy scenarios, the nine theft conditions for H3 and the six unresolved conditions for H4. We report 95\% intervals for every paired difference from cluster bootstrap resampling (10,000 resamples) and exact two-sided sign-flip tests over clusters. With $m$ clusters whose difference is not zero, the smallest attainable $p$ is $2^{1-m}$, and we mark tests that reach it. The plan's Holm correction covers the five primary tests; we also report Holm within the secondary analysis and across all ten tests.

Two dependencies in the theft design were not addressed by the plan. All pairs of a story share that story's unchanged case in each run, so one error in its checklist recurs across conditions; we report wrong changes with each repeated error counted once. Each run of the unchanged case is scored against the same run of the variant; we also report the average over all nine pairings. Analyses not in the plan are marked post hoc: two-valued execution, the strict measure, handoff causes, the routing presumption, the family-admission repair and the break-even analysis.
